\documentclass[12pt,a4paper]{article}

% Persian language support
\usepackage{xepersian}
\settextfont{IRANSansXFaNum}
\setlatintextfont{IRANSansXFaNum}
\setdigitfont{IRANSansXFaNum}

% Additional font definitions
\newfontfamily\yekanfont{Yekan}
\newfontfamily\iranfont{IRANSansXFaNum}

% Page setup and packages
\usepackage[margin=2.5cm]{geometry}
\usepackage{graphicx}
\usepackage{float}
\usepackage{amsmath}
\usepackage{booktabs}
\usepackage{longtable}
\usepackage{array}

% Line spacing
\usepackage{setspace}
\onehalfspacing

% Footnote formatting
\renewcommand{\thefootnote}{\arabic{footnote}}
% Make footnotes left-to-right for English text and position them correctly
\usepackage[bottom,flushmargin]{footmisc}
\let\oldfootnote\footnote
\renewcommand{\footnote}[1]{\oldfootnote{\lr{#1}}}

% Section numbering - use Arabic numerals with proper formatting
\renewcommand{\thesection}{\lr{\arabic{section}}}
\renewcommand{\thesubsection}{\lr{\arabic{section}.\arabic{subsection}}}
\renewcommand{\thesubsubsection}{\lr{\arabic{section}.\arabic{subsection}.\arabic{subsubsection}}}

% Fix table of contents numbering
\usepackage{tocloft}
\renewcommand{\cftsecpresnum}{\lr{}}
\renewcommand{\cftsubsecpresnum}{\lr{}}
\renewcommand{\cftsubsubsecpresnum}{\lr{}}

% Title page setup
\title{
\begin{center}
\includegraphics[height=3cm]{500px-ATU.png}\\[0.5cm]
{\large دانشگاه علامه طباطبائی}\\
{\small \lr{Allameh Tabataba'i University (ATU)}}\\[0.3cm]
{\normalsize دانشکده‌ی آمار، ریاضی و رایانه}\\
{\small \lr{Faculty of Statistics, Mathematics, and Computer}}\\[1cm]
{\yekanfont \huge گزارش پروژه}\\[0.5cm]
{\Large پیاده‌سازی سیستم پشتیبانی تصمیم‌گیری برای مراقبت‌های بهداشتی}
\end{center}
}

\author{{\iranfont \large حامد حیدریان، ایمان منتظمی، محمدمهدی عباسی}\\[0.5cm]
{\normalsize درس: سیستم‌های تصمیم‌یار}\\
{\normalsize استاد: دکتر رشیدی}\\
{\normalsize نسخه سیستم: \lr{1.0.0}}\\
{\normalsize تاریخ تکمیل: شهریور ۱۴۰۴}}

\date{}

\begin{document}

% Title page
\maketitle
\thispagestyle{empty}

% Abstract page
\newpage
\thispagestyle{empty}
\begin{center}
{\Large \textbf{چکیده}}
\end{center}

\vspace{1cm}

این پروژه به پیاده‌سازی یک سیستم پشتیبانی تصمیم‌گیری جامع برای مراقبت‌های بهداشتی می‌پردازد که بر اساس چارچوب نظری توربان-شاردا-دلن و مدل چهار مرحله‌ای هربرت سایمون طراحی شده است. سیستم از \lr{12} مجموعه‌داده مختلف بهداشتی (\lr{6} واقعی و \lr{6} مصنوعی) با بیش از \lr{24,000} رکورد استفاده می‌کند و \lr{8} الگوریتم یادگیری ماشین مختلف را پیاده‌سازی کرده است.

نتایج نشان می‌دهد که سیستم در تشخیص سرطان پستان به دقت \lr{98.2\%} و در طبقه‌بندی شراب به دقت \lr{100\%} دست یافته است. سیستم شامل ۵ داشبورد تخصصی برای نقش‌های مختلف (کادر بالینی، مدیریت اجرایی، رهبری بحرانی، مدیر بخش و مدیر مالی) است و در \lr{159} تست جامع، نرخ موفقیت \lr{89.94\%} کسب کرده است.

\textbf{کلمات کلیدی:} سیستم پشتیبانی تصمیم‌گیری، مراقبت‌های بهداشتی، یادگیری ماشین، داشبورد تخصصی، چارچوب توربان-شاردا-دلن

% Table of contents
\newpage
\tableofcontents
\thispagestyle{empty}

% List of figures
\newpage
\listoffigures
\thispagestyle{empty}

% List of tables
\newpage
\listoftables
\thispagestyle{empty}

% Start main content
\newpage
\setcounter{page}{1}

\section{مقدمه}

\subsection{تعریف مسئله تحقیق}

{\yekanfont \textbf{مسئله اصلی}}: سیستم‌های بهداشتی امروز با چالش یکپارچه‌سازی\footnote{Data Integration} داده‌های متنوع و ارائه پشتیبانی تصمیم‌گیری مناسب برای نقش‌های مختلف مواجه هستند. سیستم‌های پشتیبانی تصمیم‌گیری\footnote{Decision Support System (DSS)} موجود اغلب:

\begin{itemize}
\item قادر به پردازش همزمان\footnote{Concurrent Processing} انواع مختلف داده‌های بهداشتی نیستند
\item رابط کاربری\footnote{User Interface} یکسان برای نقش‌های مختلف ارائه می‌دهند
\item از الگوریتم‌های یادگیری ماشین\footnote{Machine Learning Algorithms} مدرن استفاده نمی‌کنند
\end{itemize}

{\iranfont \textbf{سؤال تحقیق}}: آیا می‌توان با ترکیب چارچوب نظری\footnote{Theoretical Framework} سیستم‌های پشتیبانی تصمیم‌گیری، الگوریتم‌های یادگیری ماشین مدرن، و طراحی مبتنی بر نقش\footnote{Role-based Design}، سیستمی طراحی کرد که بتواند مسائل مختلف بهداشتی را به‌طور مؤثر حل کند؟

\textbf{فرضیه تحقیق}: پیاده‌سازی سیستم پشتیبانی تصمیم‌گیری مبتنی بر چارچوب توربان-شاردا-دلن\footnote{Turban-Sharda-Delen Framework} با استفاده از الگوریتم‌های یادگیری ماشین و طراحی مبتنی بر نقش، می‌تواند عملکرد تصمیم‌گیری در حوزه‌های مختلف بهداشتی را بهبود دهد.

کتاب «تحلیل‌ها، علم داده و هوش مصنوعی: سیستم‌های پشتیبانی تصمیم‌گیری»\footnote{Analytics, Data Science, \& Artificial Intelligence: Systems for Decision Support (11th Edition)} نوشته توربان، شاردا و دلن چارچوب جامعی برای طراحی و پیاده‌سازی سیستم‌های پشتیبانی تصمیم‌گیری ارائه می‌دهد. 

این چارچوب بر اساس مدل چهار مرحله‌ای هربرت سایمون\footnote{Herbert Simon's Four-Phase Decision Model} شامل مراحل هوشمندی\footnote{Intelligence Phase}، طراحی\footnote{Design Phase}، انتخاب\footnote{Choice Phase} و پیاده‌سازی\footnote{Implementation Phase} است.

در این پروژه، یک سیستم پشتیبانی تصمیم‌گیری جامع برای مراقبت‌های بهداشتی\footnote{Healthcare} پیاده‌سازی شده است که از چهار زیرسیستم اصلی تشکیل شده: مدیریت داده‌ها\footnote{Data Management Subsystem}، مدیریت مدل‌ها\footnote{Model Management Subsystem}، مدیریت دانش\footnote{Knowledge Management Subsystem} و رابط کاربری. سیستم با استفاده از الگوریتم‌های یادگیری ماشین مختلف شامل جنگل تصادفی\footnote{Random Forest}، تقویت گرادیان شدید\footnote{XGBoost (Extreme Gradient Boosting)}، تقویت گرادیان سبک\footnote{LightGBM (Light Gradient Boosting Machine)} و ماشین بردار پشتیبان\footnote{Support Vector Machine (SVM)}، قابلیت تحلیل و پیش‌بینی در حوزه‌های مختلف بهداشتی را فراهم می‌آورد.

\subsection{مسائل حل شده توسط سیستم}

سیستم طراحی شده برای حل چهار دسته مسئله اصلی در مراقبت‌های بهداشتی طراحی شده است:

\textbf{۱. مسائل تشخیص و پیش‌بینی بالینی}: شامل تشخیص زودهنگام\footnote{Early Detection} سرطان پستان\footnote{Breast Cancer}، پیش‌بینی پیشرفت دیابت\footnote{Diabetes Progression Prediction}، و ارزیابی ریسک\footnote{Risk Assessment} بستری مجدد بیماران که از طریق داشبورد کادر بالینی\footnote{Clinical Staff Dashboard} قابل دسترسی است.

\textbf{۲. مسائل مدیریت منابع و ظرفیت}: شامل بهینه‌سازی ظرفیت\footnote{Capacity Optimization} بیمارستان، تخصیص بهینه\footnote{Optimal Allocation} تخت‌ها، و مدیریت جریان بیماران\footnote{Patient Flow Management} که از طریق داشبورد مدیریت اجرایی\footnote{Administrative Executive Dashboard} پشتیبانی می‌شود.

\textbf{۳. مسائل ارزیابی عملکرد}: شامل نظارت بر عملکرد کادر\footnote{Staff Performance Monitoring}، ارزیابی کیفیت بخش‌ها\footnote{Department Quality Assessment}، و تحلیل رضایت بیماران\footnote{Patient Satisfaction Analysis} که از طریق داشبورد مدیر بخش\footnote{Department Manager Dashboard} قابل پیگیری است.

\textbf{۴. مسائل مالی و استراتژیک}: شامل تحلیل هزینه‌های بهداشتی\footnote{Healthcare Cost Analysis}، پیش‌بینی درآمد\footnote{Revenue Prediction}، و بهینه‌سازی بودجه\footnote{Budget Optimization} که از طریق داشبورد مدیر مالی\footnote{Financial Manager Dashboard} و رهبری بحرانی\footnote{Critical Leadership Dashboard} پشتیبانی می‌شود.

\section{روش کار}

این پروژه با هدف حل مسائل عملی در سیستم‌های بهداشتی، از روش‌شناسی ترکیبی استفاده کرده که شامل سه رویکرد اصلی است:

\subsection{روش‌شناسی کلی پروژه}

{\yekanfont \textbf{رویکرد اول - چارچوب نظری}}: پیاده‌سازی بر اساس چارچوب نظری کتاب توربان، شاردا و دلن و مدل چهار مرحله‌ای هربرت سایمون برای تصمیم‌گیری انجام شده است.

{\iranfont \textbf{رویکرد دوم - روش‌شناسی کریسپ-دی‌ام}}: برای توسعه مدل‌های یادگیری ماشین از روش‌شناسی شش مرحله‌ای کریسپ-دی‌ام\footnote{Cross-Industry Standard Process for Data Mining (CRISP-DM)} استفاده شده که شامل درک کسب‌وکار، درک داده‌ها، آماده‌سازی داده‌ها، مدل‌سازی، ارزیابی و استقرار است.

\textbf{رویکرد سوم - آزمایش تجربی}: برای اعتبارسنجی فرضیه‌های تحقیق، طرح آزمایشی جامع طراحی شده که شامل:
\begin{itemize}
\item آزمایش عملکرد ۸ الگوریتم یادگیری ماشین بر روی ۱۲ مجموعه‌داده مختلف
\item مقایسه تطبیقی الگوریتم‌ها با معیارهای استاندارد (دقت\footnote{Accuracy}، دقت مثبت\footnote{Precision}، حساسیت\footnote{Recall/Sensitivity}، امتیاز F1\footnote{F1-Score})
\item ارزیابی کیفیت سیستم از طریق ۱۵۹ تست واحد و یکپارچگی
\item تحلیل آماری نتایج برای تعیین بهترین الگوریتم برای هر نوع مسئله
\end{itemize}

\subsection{داده‌ها و آماده‌سازی}

\subsubsection{مجموعه‌داده‌های مورد استفاده}

سیستم از ۱۲ مجموعه‌داده مختلف استفاده می‌کند که شامل داده‌های واقعی\footnote{Real-world Data} و مصنوعی\footnote{Synthetic Data} است:

{\yekanfont \textbf{مجموعه‌داده‌های واقعی}}:
\begin{itemize}
\item \textbf{دیابت}\footnote{Diabetes Dataset} (\lr{442} نمونه، \lr{10} ویژگی): پیش‌بینی پیشرفت بیماری دیابت با متغیر هدف رگرسیون\footnote{Regression Task Type}
\item \textbf{سرطان پستان}\footnote{Breast Cancer Dataset} (\lr{569} نمونه، \lr{30} ویژگی): تشخیص سرطان پستان با دقت بالا
\item \textbf{شراب}\footnote{Wine Dataset} (\lr{178} نمونه، \lr{13} ویژگی): طبقه‌بندی انواع شراب برای تحقیقات پزشکی
\item \textbf{ظرفیت بیمارستان}\footnote{Hospital Capacity Dataset} (\lr{20,640} نمونه، \lr{8} ویژگی): داده‌های ظرفیت بیمارستان
\item \textbf{هزینه‌های بهداشتی جهانی}\footnote{Global Healthcare Expenditure Dataset} (\lr{117} نمونه، \lr{14} ویژگی): هزینه‌های بهداشتی \lr{16} کشور از \lr{2015} تا \lr{2024}
\end{itemize}

{\iranfont \textbf{مجموعه‌داده‌های مصنوعی}}:
\begin{itemize}
\item \textbf{جمعیت‌شناسی بیماران}\footnote{Patient Demographics Synthetic Dataset} (\lr{1,000} نمونه، \lr{20} ویژگی)
\item \textbf{نتایج بالینی}\footnote{Clinical Outcomes Synthetic Dataset} (\lr{1,000} نمونه، \lr{7} ویژگی)
\item \textbf{عملکرد کادر}\footnote{Staff Performance Synthetic Dataset} (\lr{200} نمونه، \lr{9} ویژگی)
\item \textbf{عملکرد بخش‌ها}\footnote{Department Performance Synthetic Dataset} (\lr{7} نمونه، \lr{8} ویژگی)
\item \textbf{معیارهای مالی}\footnote{Financial Metrics Synthetic Dataset} (\lr{24} نمونه، \lr{8} ویژگی)
\item \textbf{اثربخشی دارو}\footnote{Medication Effectiveness Dataset} (\lr{178} نمونه، \lr{14} ویژگی)
\end{itemize}

\subsection{پیش‌پردازش داده‌ها}

سیستم از کلاس موتور پیش‌پردازش\footnote{Preprocessing Engine} برای پیش‌پردازش هوشمند داده‌ها استفاده می‌کند که شامل:

\textbf{مراحل پیش‌پردازش برای طبقه‌بندی}:
\begin{enumerate}
\item مدیریت مقادیر گمشده\footnote{Missing Values Management} با استراتژی میانه\footnote{Median Strategy}
\item کدگذاری متغیرهای کیفی با روش کدگذاری برچسب\footnote{Label Encoding Method}
\item نرمال‌سازی ویژگی‌های عددی با مقیاس‌کننده استاندارد\footnote{Standard Scaler}
\item متعادل‌سازی مجموعه‌داده با روش اسموت\footnote{Synthetic Minority Oversampling Technique (SMOTE)}
\item انتخاب ویژگی با روش اطلاعات متقابل\footnote{Mutual Information Feature Selection}
\end{enumerate}

\subsection{طراحی رابط کاربری مبتنی بر نقش}

برای حل مسائل مختلف سیستم‌های بهداشتی، از روش طراحی مبتنی بر نقش\footnote{Role-based Design Approach} استفاده شده است. پنج داشبورد تخصصی طراحی و پیاده‌سازی شده است:

\begin{figure}[H]
\centering
\includegraphics[width=0.8\textwidth]{image_Clinical_Staff.png}
\caption{داشبورد کادر بالینی - پیاده‌سازی مدل‌های تشخیص بالینی}
\label{fig:clinical_dashboard}
\end{figure}

\begin{figure}[H]
\centering
\includegraphics[width=0.8\textwidth]{image_Administrative_Executive.png}
\caption{داشبورد مدیریت اجرایی - سیستم نظارت بر عملکرد و منابع}
\label{fig:admin_dashboard}
\end{figure}

\begin{figure}[H]
\centering
\includegraphics[width=0.8\textwidth]{image_Critical_Leadership.png}
\caption{داشبورد رهبری بحرانی - تحلیل استراتژیک و پیش‌بینی روندها}
\label{fig:leadership_dashboard}
\end{figure}

\begin{figure}[H]
\centering
\includegraphics[width=0.8\textwidth]{image_Department_Manager.png}
\caption{داشبورد مدیر بخش - سیستم مدیریت کیفیت و عملکرد بخش}
\label{fig:department_dashboard}
\end{figure}

\begin{figure}[H]
\centering
\includegraphics[width=0.8\textwidth]{image_Financial_Manager.png}
\caption{داشبورد مدیر مالی - سیستم تحلیل مالی و کنترل بودجه}
\label{fig:financial_dashboard}
\end{figure}

\subsection{پیاده‌سازی مدل}

\subsubsection{نوع مدل و ساختار شبکه}

سیستم از معماری مدولار\footnote{Modular Architecture} استفاده می‌کند که شامل چهار زیرسیستم اصلی است:

{\yekanfont \textbf{الگوریتم‌های طبقه‌بندی}}:
\begin{itemize}
\item طبقه‌بند جنگل تصادفی: برای طبقه‌بندی دوکلاسه و چندکلاسه\footnote{Binary and Multi-class Classification}
\item رگرسیون لجستیک\footnote{Logistic Regression}: برای طبقه‌بندی خطی\footnote{Linear Classification}
\item درخت تصمیم\footnote{Decision Tree}: برای طبقه‌بندی قابل تفسیر\footnote{Interpretable Classification}
\item طبقه‌بند ایکس‌جی‌بوست\footnote{XGBoost Classifier}: برای عملکرد بالا\footnote{High Performance}
\item ماشین بردار پشتیبان: برای طبقه‌بندی با هسته‌های مختلف\footnote{Different Kernel Classification}
\end{itemize}

{\iranfont \textbf{الگوریتم‌های رگرسیون}}:
\begin{itemize}
\item رگرسور جنگل تصادفی: برای پیش‌بینی مقادیر پیوسته\footnote{Continuous Value Prediction}
\item رگرسیون خطی\footnote{Linear Regression}: برای رگرسیون خطی\footnote{Linear Regression Method}
\item رگرسور ایکس‌جی‌بوست\footnote{XGBoost Regressor}: برای رگرسیون با دقت بالا\footnote{High Accuracy Regression}
\item رگرسیون بردار پشتیبان\footnote{Support Vector Regression}: برای رگرسیون با ماشین بردار پشتیبان\footnote{SVM-based Regression}
\end{itemize}

\section{ارزیابی و نتایج}

\subsection{روش و معیار ارزیابی}

ارزیابی سیستم در دو بعد اصلی انجام شده است:

\textbf{ارزیابی فنی مدل‌ها}:
\begin{itemize}
\item معیارهای طبقه‌بندی: دقت، دقت مثبت، حساسیت، امتیاز اف‌یک، سطح زیر منحنی راک\footnote{Area Under ROC Curve}
\item اعتبارسنجی متقابل: اعتبارسنجی متقابل پنج‌تایی\footnote{5-Fold Cross-Validation}
\item مقایسه مدل‌ها: مقایسه عملکرد الگوریتم‌های مختلف
\end{itemize}

\textbf{ارزیابی سیستم}:
\begin{itemize}
\item تست‌های واحد: ۱۵۹ تست جامع برای بررسی عملکرد ماژول‌ها
\item تست‌های یکپارچگی: بررسی عملکرد کل سیستم
\item ارزیابی کیفیت کد: بررسی استانداردهای کدنویسی
\item تست عملکرد: سنجش سرعت و کارایی سیستم
\end{itemize}

\subsection{آزمایش‌ها و نتایج}

\subsubsection{نتایج مدل‌های یادگیری ماشین}

بر اساس آزمایش‌های انجام شده، نتایج دقیق به شرح زیر است:

{\yekanfont \textbf{مجموعه‌داده دیابت}} (\lr{442} نمونه، \lr{10} ویژگی، رگرسیون):
\begin{itemize}
\item بهترین مدل: رگرسیون خطی با \lr{$R^2 = 0.453$}
\item جنگل تصادفی: \lr{$R^2 = 0.443$}، \lr{$RMSE = 54.332$}
\item درخت تصمیم: \lr{$R^2 = 0.061$}، \lr{$RMSE = 70.546$}
\end{itemize}

{\iranfont \textbf{مجموعه‌داده سرطان پستان}} (\lr{569} نمونه، \lr{30} ویژگی، طبقه‌بندی):
\begin{itemize}
\item بهترین مدل: رگرسیون لجستیک با دقت \lr{98.2\%}
\item جنگل تصادفی: دقت = \lr{95.6\%}، \lr{$F_1 = 0.956$}
\item درخت تصمیم: دقت = \lr{91.2\%}، \lr{$F_1 = 0.913$}
\end{itemize}

\textbf{مجموعه‌داده شراب} (\lr{178} نمونه، \lr{13} ویژگی، طبقه‌بندی):
\begin{itemize}
\item بهترین مدل: جنگل تصادفی با دقت \lr{100\%}
\item رگرسیون لجستیک: دقت = \lr{97.2\%}، \lr{$F_1 = 0.972$}
\item درخت تصمیم: دقت = \lr{94.4\%}، \lr{$F_1 = 0.945$}
\end{itemize}

\subsubsection{نتایج تست‌های سیستم}

\textbf{عملکرد کلی سیستم}:
\begin{itemize}
\item تعداد کل تست‌ها: \lr{159} تست
\item نرخ موفقیت کلی: \lr{89.94\%}
\item زمان اجرای کل: \lr{15.86} ثانیه
\item تعداد مجموعه تست: ۶ مجموعه
\end{itemize}

\begin{table}[H]
\centering
\caption{تفکیک عملکرد ماژول‌ها}
\label{tab:module_performance}
\begin{tabular}{>{\centering\arraybackslash}p{3cm}>{\centering\arraybackslash}p{3.5cm}>{\centering\arraybackslash}p{2cm}>{\centering\arraybackslash}p{2.5cm}>{\centering\arraybackslash}p{2.5cm}}
\toprule
\textbf{مجموعه تست} & \textbf{ماژول} & \textbf{تعداد تست} & \textbf{نرخ موفقیت} & \textbf{زمان اجرا} \\
\midrule
\lr{TestSuite\_1} & مدیریت مدل‌ها & \lr{25} & \lr{60.0\%} & \lr{1.03} ثانیه \\
\lr{TestSuite\_2} & مدیریت داده‌ها & \lr{29} & \lr{100\%} & \lr{0.70} ثانیه \\
\lr{TestSuite\_3} & ماژول‌های رابط کاربری & \lr{28} & \lr{78.57\%} & \lr{1.75} ثانیه \\
\lr{TestSuite\_4} & تحلیل‌ها & \lr{31} & \lr{100\%} & \lr{1.11} ثانیه \\
\lr{TestSuite\_5} & مدیریت دانش & \lr{11} & \lr{100\%} & \lr{2.31} ثانیه \\
\lr{TestSuite\_6} & اجزای رابط کاربری & \lr{35} & \lr{100\%} & \lr{3.31} ثانیه \\
\bottomrule
\end{tabular}
\end{table}

\section{جمع‌بندی و پیشنهادها}

\subsection{جمع‌بندی کلی از نتایج}

پروژه سیستم پشتیبانی تصمیم‌گیری مراقبت‌های بهداشتی با موفقیت قابل‌توجه پیاده‌سازی شده و نتایج ملموسی حاصل کرده است. سیستم با ترکیب چهار زیرسیستم اصلی، پلتفرمی جامع و کاربردی برای پشتیبانی تصمیمات بهداشتی ارائه داده است.

{\yekanfont \textbf{دستاوردهای کلیدی}}:
\begin{itemize}
\item پیاده‌سازی کامل سیستم: \lr{12} مجموعه‌داده مختلف بهداشتی با \lr{24,000+} رکورد
\item عملکرد بالای مدل‌ها: دقت \lr{98.2\%} در تشخیص سرطان پستان، \lr{100\%} در طبقه‌بندی شراب
\item معماری مدولار: چهار زیرسیستم اصلی با \lr{159} تست جامع (\lr{89.94\%} موفقیت)
\item داشبوردهای تخصصی\footnote{Specialized Dashboards}: ۵ نوع داشبورد برای نقش‌های مختلف
\item تنوع الگوریتم‌ها: \lr{8} الگوریتم یادگیری ماشین با \lr{24} مدل آموزش‌دیده
\end{itemize}

\subsection{چالش‌های موجود در روند کار}

{\iranfont \textbf{چالش‌های فنی شناسایی شده}}:
\begin{itemize}
\item مسائل ماژول مدیریت مدل‌ها (۶۰٪ موفقیت)
\item مسائل رابط کاربری (۷۸.۵۷٪ موفقیت)
\item چالش‌های معماری در مدیریت گواهی‌های امنیتی
\end{itemize}

\textbf{چالش‌های داده‌ای}:
\begin{itemize}
\item وجود \lr{68} مقدار گمشده در مجموعه‌داده هزینه‌های بهداشتی
\item عدم تعادل در اندازه مجموعه‌داده‌ها (از \lr{7} تا \lr{20,640} نمونه)
\end{itemize}

\subsection{پیشنهادها برای بهبود پروژه و مسیرهای آتی}

\textbf{بهبودهای فوری (اولویت بالا)}:
\begin{itemize}
\item تکمیل مدیر مدل: پیاده‌سازی متدهای گمشده
\item تکمیل مدیر داشبورد: افزودن متدهای ایجاد داشبورد
\item رفع خطاهای رابط کاربری
\item بهبود پوشش تست: افزایش نرخ موفقیت از \lr{89.94\%} به \lr{95\%+}
\end{itemize}

\textbf{توسعه‌های میان‌مدت}:
\begin{itemize}
\item افزایش قابلیت‌ها: اضافه کردن مجموعه‌داده‌های بیشتر
\item پیاده‌سازی الگوریتم‌های یادگیری عمیق\footnote{Deep Learning Algorithms Implementation}
\item بهبود رابط کاربری: طراحی داشبوردهای تعاملی‌تر\footnote{More Interactive Dashboard Design}
\item یکپارچگی: اتصال به پایگاه‌داده‌های بیمارستانی واقعی\footnote{Real Hospital Database Connection}
\end{itemize}

\textbf{چشم‌انداز بلندمدت}:
\begin{itemize}
\item فناوری‌های نوین: پیاده‌سازی مدل‌های زبانی بزرگ\footnote{Large Language Models Implementation}
\item استفاده از بینایی کامپیوتری\footnote{Computer Vision Usage} برای تحلیل تصاویر پزشکی\footnote{Medical Image Analysis}
\item یکپارچگی با اینترنت اشیا\footnote{Internet of Things Integration} و سنسورهای پزشکی\footnote{Medical Sensors}
\end{itemize}

\section{منابع و مراجع}

\begin{enumerate}
\item توربان، ای.، شاردا، آر.، و دلن، د. (۲۰۲۰). \textit{تحلیل‌ها، علم داده و هوش مصنوعی: سیستم‌های پشتیبانی تصمیم‌گیری} (ویرایش یازدهم). پیرسون.

\item سایمون، هربرت ای. (۱۹۶۰). \textit{علم جدید تصمیم‌گیری مدیریت}. هارپر و رو.

\item فرایند استاندارد صنعتی برای داده‌کاوی (کریسپ-دی‌ام) (۲۰۰۰).

\item کتابخانه‌های یادگیری ماشین:
\begin{itemize}
\item ساینت‌کیت لرن (>=۱.۳.۰): https://scikit-learn.org/
\item ایکس‌جی‌بوست (>=۱.۷.۰): https://xgboost.readthedocs.io/
\item لایت‌جی‌بی‌ام (>=۴.۰.۰): https://lightgbm.readthedocs.io/
\end{itemize}

\item هوش مصنوعی قابل تفسیر:
\begin{itemize}
\item شپ (>=۰.۴۲.۰): https://shap.readthedocs.io/
\item لایم (>=۰.۲.۰): https://lime-ml.readthedocs.io/
\end{itemize}

\item داشبورد و تجسم داده‌ها:
\begin{itemize}
\item استریم‌لیت (>=۱.۲۸.۰): https://streamlit.io/
\item پلاتلی (>=۵.۱۵.۰): https://plotly.com/
\end{itemize}

\item مخزن یادگیری ماشین دانشگاه کالیفرنیا: \\
\lr{https://archive.ics.uci.edu/ml/}

\item پایگاه داده هزینه‌های بهداشتی جهانی سازمان بهداشت جهانی: \\
\lr{https://apps.who.int/nha/database}
\end{enumerate}

\end{document}